\documentclass[aspectratio=169,spanish,xcolor=table]{beamer}
\usepackage{iftex}
\ifPDFTeX
  \usepackage[utf8]{inputenc}
  \usepackage[T1]{fontenc}
\fi
\usepackage[spanish,es-tabla,es-noshorthands]{babel}
\usepackage{helvet}
\renewcommand{\familydefault}{\sfdefault}
\usepackage{tikz,pgfplots,array,booktabs,ragged2e}
\pgfplotsset{compat=1.18}
\usetikzlibrary{arrows.meta,positioning,calc,shapes.geometric,fit}

\definecolor{iiiepeNavy}{HTML}{163B5C}
\definecolor{iiiepeBlue}{HTML}{1F6F9F}
\definecolor{iiiepeAqua}{HTML}{20B7A8}
\definecolor{iiiepeOrange}{HTML}{F28C28}
\definecolor{iiiepePaper}{HTML}{F6F8F7}
\definecolor{iiiepeInk}{HTML}{1E2A32}
\definecolor{iiiepePale}{HTML}{EEF5F8}
\definecolor{iiiepeMint}{HTML}{E7F6F4}
\definecolor{iiiepeWarm}{HTML}{FFF3E6}

\usetheme{default}
\usecolortheme{default}
\setbeamertemplate{navigation symbols}{}
\setbeamercolor{normal text}{fg=iiiepeInk,bg=white}
\setbeamercolor{structure}{fg=iiiepeNavy}
\setbeamercolor{frametitle}{fg=white,bg=iiiepeNavy}
\setbeamercolor{block title}{fg=white,bg=iiiepeBlue}
\setbeamercolor{block body}{fg=iiiepeInk,bg=iiiepePale}
\setbeamercolor{item}{fg=iiiepeOrange}

\setbeamertemplate{frametitle}{%
  \nointerlineskip%
  \begin{beamercolorbox}[wd=\paperwidth,ht=0.95cm,dp=0.22cm,leftskip=0.45cm,rightskip=0.45cm]{frametitle}%
    \begin{minipage}[c]{13.2cm}\usebeamerfont{frametitle}\large\bfseries\insertframetitle\end{minipage}%
    \hfill\begin{minipage}[c]{1.6cm}\raggedleft\includegraphics[height=0.56cm]{img/iiiepe/logo-iiiepe-blanco.png}\end{minipage}%
  \end{beamercolorbox}%
}

\setbeamertemplate{footline}{%
  \leavevmode%
  \hbox{%
    \begin{beamercolorbox}[wd=.78\paperwidth,ht=0.5cm,dp=0.18cm,leftskip=0.45cm]{frametitle}%
      \scriptsize\color{white}M. de la Cruz Muñoz · MAEM · Matrícula 00874 · SD26 · Actividad 6 · Tercer Momento%
    \end{beamercolorbox}%
    \begin{beamercolorbox}[wd=.22\paperwidth,ht=0.5cm,dp=0.18cm,rightskip=0.45cm]{frametitle}%
      \hfill\scriptsize\color{white}\insertframenumber{} / \inserttotalframenumber%
    \end{beamercolorbox}%
  }%
  \vskip0pt%
}

\newcommand{\Pill}[2]{%
  \tikz[baseline=(x.base)]{\node[rounded corners=3pt,fill=#1!14,draw=#1!75!black,thick,inner xsep=6pt,inner ysep=4pt,align=center] (x) {\scriptsize\bfseries #2};}%
}

\title{Tres Momentos: Diagnóstico e Innovación Matemática}
\subtitle{Actividad 6 · Bloque III · Módulo 51809}
\author{Mtro. Martín Jonathan de la Cruz Muñoz · Matrícula 00874}
\institute{IIIEPE · Maestría en Aprendizaje y Enseñanza de las Matemáticas · Grupo A}
\date{19 de noviembre de 2026}

\begin{document}

\begin{frame}[plain]
  \begin{tikzpicture}[remember picture,overlay]
    \fill[iiiepeNavy] (current page.north west) rectangle ([xshift=10.3cm]current page.south west);
    \fill[iiiepeBlue] ([xshift=10.3cm]current page.north west) rectangle (current page.south east);
    \fill[iiiepeOrange] ([xshift=0.8cm,yshift=-6.45cm]current page.north west) rectangle ([xshift=9.3cm,yshift=-6.53cm]current page.north west);
    \node[anchor=north west,text width=8.7cm,align=left] at ([xshift=0.8cm,yshift=-1.15cm]current page.north west) {%
      {\color{iiiepeAqua!45!white}\small\bfseries MAESTRÍA EN APRENDIZAJE Y ENSEÑANZA DE LAS MATEMÁTICAS · SD26\par}%
      \vspace{0.35cm}
      {\color{white}\Large\bfseries \inserttitle\par}%
      \vspace{0.25cm}
      {\color{white!90}\normalsize \insertsubtitle\par}%
    };
    \node[anchor=south west,text width=8.7cm,align=left] at ([xshift=0.8cm,yshift=0.75cm]current page.south west) {%
      {\color{white}\normalsize\bfseries Mtro. Martín Jonathan de la Cruz Muñoz\par}%
      {\color{white!85}\scriptsize Matrícula 00874 · martin.delacruz@campus.iiiepe.edu.mx\par}%
      {\color{iiiepeAqua!35!white}\scriptsize Monterrey, Nuevo León · 19 de noviembre de 2026\par}%
    };
    \node[anchor=center,text width=5.0cm,align=center] at ([xshift=13.15cm,yshift=-4.5cm]current page.north west) {%
      \includegraphics[width=3.6cm]{img/iiiepe/logo-iiiepe-blanco.png}\par\vspace{0.45cm}%
      {\color{white}\large\bfseries IIIEPE\par}%
      \vspace{0.1cm}%
      {\color{white!85}\scriptsize Fundamentos para la enseñanza\\y el aprendizaje I\par}%
    };
  \end{tikzpicture}
\end{frame}
\begin{frame}{Ruta del Bloque III: Tres Momentos Integrados}
  \centering
  \Pill{iiiepeNavy}{1. Síntesis crítica (Quiroga y Lara, 2022)} \hspace{0.2cm}
  \Pill{iiiepeBlue}{2. Diagnóstico en 6 aspectos} \hspace{0.2cm}
  \Pill{iiiepeOrange!85!black}{3. Innovación en 7 componentes}
  \vspace{0.45cm}
  \begin{block}{Tesis del ejercicio (Módulo 51809)}
    Superar la mecanización algebraica sin sentido físico exige revisar nuestras decisiones didácticas antes de culpar al estudiantado y diseñar retos matemáticos orientados al aprendizaje profundo \textbf{(propuesta por realizar)}.
  \end{block}
\end{frame}

\begin{frame}{Primer Momento: Cinco Ejes de Quiroga Aguilar y Lara Órdenes (2022)}
  \begin{columns}[T]
    \column{0.48\textwidth}
      \begin{itemize}\scriptsize
        \item \textbf{1. Superficial vs Profundo:} Reproducción pasiva frente a comprensión y transferencia \textbf{(Marton y Säljö, 1976)}.
        \item \textbf{2. Tres dimensiones:} Integración de cognitivo, interpersonal e intrapersonal.
        \item \textbf{3. Cambio de roles:} Del control docente hacia asociación y autonomía \textbf{(Fullan y Langworthy, 2014)}.
      \end{itemize}
    \column{0.48\textwidth}
      \begin{itemize}\scriptsize
        \item \textbf{4. Nuevas pedagogías:} Problemas auténticos del entorno, co-construcción y liderazgo pedagógico \textbf{(Rivas et al., 2017)}.
        \item \textbf{5. Creencias y cultura escolar:} No basta renovar recursos o actividades; hay que transformar creencias compartidas sobre qué significa aprender matemáticas.
      \end{itemize}
  \end{columns}
\end{frame}

\begin{frame}{Segundo Momento: Diagnóstico Situado (Aspectos 1--3)}
  \scriptsize\renewcommand{\arraystretch}{1.35}
  \begin{tabular}{p{2.7cm}p{10.3cm}}
    \toprule
    \rowcolor{iiiepePale} \textbf{Aspecto} & \textbf{Hallazgo observado en evidencias de aula} \\
    \midrule
    1. Problemática & Revisión de cuadernos muestra cálculo automático de $m$ y despejes sin explicar significado físico, unidades ni validez en modelos lineales. \\
    2. Práctica actual & Tendencia a iniciar con fórmula cerrada en pizarrón, ejemplos guiados paso a paso y ejercitación individual aislada. \\
    3. Concepciones & Supuesto implícito de que dominar el algoritmo previo precede a la interpretación física del problema. \\
    \bottomrule
  \end{tabular}
\end{frame}

\begin{frame}{Segundo Momento: Diagnóstico Situado (Aspectos 4--6)}
  \scriptsize\renewcommand{\arraystretch}{1.35}
  \begin{tabular}{p{2.7cm}p{10.3cm}}
    \toprule
    \rowcolor{iiiepePale} \textbf{Aspecto} & \textbf{Hallazgo y reflexión crítica sin culpar solo al alumnado} \\
    \midrule
    4. Aprendizaje & Evidencias de enfoque superficial ante cambios de notación; interés profundo incipiente cuando se conectan gráficas con sensores de taller. \\
    5. Causas posibles & Exceso de tareas repetitivas descontextualizadas, prisa por cerrar el tema y evaluación centrada en resultado numérico final. \\
    6. Necesidad de cambio & Invertir el diseño: partir de desafío auténtico, habilitar trabajo colaborativo argumentado y evaluar con devolución formativa continua. \\
    \bottomrule
  \end{tabular}
\end{frame}

\begin{frame}{Tercer Momento: Propuesta «Calibrar y Comparar para Decidir»}
  \begin{columns}[T]
    \column{0.48\textwidth}
      \begin{block}{Elementos 1--3 (Contenido, Problema, Propósito)}
        \begin{itemize}\scriptsize
          \item \textbf{Contenido:} Modelación lineal ($y=mx+b$ vs $y=kx$), pendiente, sesgo inicial e intervalo de validez en bachillerato técnico (4.º sem.).
          \item \textbf{Problemática:} Algoritmia sin sentido dimensional.
          \item \textbf{Propósito:} Articular tabla, gráfica y ecuación para justificar decisiones técnicas.
        \end{itemize}
      \end{block}
    \column{0.48\textwidth}
      \begin{block}{Elemento 4: Desafío Matemático Simulado}
        \scriptsize Dos sensores de banco ($S_1$ patrón vs $S_2$: $y=1.05x+0.2$) y dos esquemas de costo ($C_A=250+80h$, $C_B=100+110h$):
        \begin{itemize}\tiny
          \item ¿Qué error introduce suponer proporcionalidad directa?
          \item ¿Desde qué umbral conviene cambiar de opción?
          \item Justificar por escrito predicción, unidades y límites.
        \end{itemize}
      \end{block}
  \end{columns}
\end{frame}

\begin{frame}{Elemento 5: Estrategias de Enseñanza (Propuesta por Realizar)}
  \centering
  \begin{tikzpicture}[font=\scriptsize, >=Latex, scale=0.85, transform shape]
    \node[draw=iiiepeNavy,fill=iiiepePale,thick,rounded corners=3pt,text width=2.7cm,minimum height=1.7cm,align=center] (a) {\textbf{Sesión 1}\\Predicción individual y tablas reales};
    \node[draw=iiiepeBlue,fill=iiiepePale,thick,rounded corners=3pt,text width=2.7cm,minimum height=1.7cm,align=center,right=0.45cm of a] (b) {\textbf{Sesión 2}\\Tríos rotativos (gráfica, unidades, vocero)};
    \node[draw=iiiepeAqua!80!black,fill=iiiepeMint,thick,rounded corners=3pt,text width=2.7cm,minimum height=1.7cm,align=center,right=0.45cm of b] (c) {\textbf{Sesión 3}\\Plenaria dialógica y crítica de errores};
    \node[draw=iiiepeOrange!90!black,fill=iiiepeWarm,thick,rounded corners=3pt,text width=2.7cm,minimum height=1.7cm,align=center,right=0.45cm of c] (d) {\textbf{Cierre}\\Segunda versión + metacognición};
    \draw[->,very thick,iiiepeNavy] (a) -- (b);
    \draw[->,very thick,iiiepeBlue] (b) -- (c);
    \draw[->,very thick,iiiepeAqua!80!black] (c) -- (d);
  \end{tikzpicture}
  \vspace{0.35cm}
  \begin{itemize}\scriptsize
    \item Participación activa desde conjeturas previas antes de despejar símbolos \textbf{(Sawyer, 2008)}.
    \item Retirada progresiva del andamiaje docente para consolidar autonomía estudiantil.
  \end{itemize}
\end{frame}

\begin{frame}{Elemento 6: Integración de las Tres Dimensiones}
  \begin{itemize}\normalsize
    \item \textbf{\color{iiiepeNavy}Cognitiva:} Traducción bidireccional tabla--gráfica--símbolo, interpretación física de pendiente y sesgo, verificación de dominio válido.
    \item \textbf{\color{iiiepeBlue}Interpersonal:} Negociación respetuosa en roles rotativos y defensa argumentada entre equipos durante la plenaria.
    \item \textbf{\color{iiiepeAqua!75!black}Intrapersonal:} Diario breve de revisión sobre supuestos iniciales corregidos, cultivando mentalidad de crecimiento frente al error.
  \end{itemize}
\end{frame}

\begin{frame}{Elemento 7: Evaluación Formativa e Indicadores}
  \scriptsize\renewcommand{\arraystretch}{1.38}
  \begin{tabular}{p{2.6cm}p{4.8cm}p{5.5cm}}
    \toprule
    \rowcolor{iiiepePale} \textbf{Eje} & \textbf{Evidencia recogida} & \textbf{Indicador de avance} \\
    \midrule
    Modelación & Registro tabular--gráfico y modelo lineal. & Diferencia variación proporcional vs afín con unidades correctas. \\
    Argumentación & Dictamen técnico comparando rangos. & Justifica decisiones señalando supuestos y límites físicos. \\
    Metacognición & Segunda versión tras devolución oral/escrita. & Corrige errores conceptuales identificando la causa del cambio. \\
    \bottomrule
  \end{tabular}
\end{frame}

\begin{frame}{Viabilidad e Impacto Esperado en el Plantel}
  \begin{block}{Carácter de la propuesta: diseño por implementar}
    \small No se reportan calificaciones ficticias ni resultados observados aún; la secuencia queda lista para ejecutarse bajo seguimiento de investigación-acción en grupos de CONALEP, CECyTE o ingeniería inicial \textbf{(Restrepo Gómez, 2004)}.
  \end{block}
  \vspace{0.3cm}
  \begin{itemize}\small
    \item Fortalece el tránsito desde una organización puramente técnica hacia el diálogo humanista y comprensivo.
    \item Promueve equidad al ofrecer roles claros y devoluciones formativas antes de evaluar sumativamente.
  \end{itemize}
\end{frame}

\begin{frame}{Conclusiones de la Actividad 6}
  \begin{itemize}\normalsize
    \item La innovación didáctica surge de problematizar con evidencia nuestra propia enseñanza, no de culpar al alumno.
    \item Integrar las tres dimensiones del aprendizaje profundo transforma funciones lineales en herramientas reales de decisión.
    \item Las evidencias y criterios formativos definidos permiten evaluar rigor y equidad al llevar el diseño al aula.
  \end{itemize}
\end{frame}

\begin{frame}{Referencias Principales (APA 7)}
  \scriptsize\setlength{\parskip}{5pt}
  \begin{itemize}
    \item Boström, E., y Palm, T. (2023). Formative assessment practice in mathematics. \textit{Frontiers in Education}, 8, 1101192.
    \item Collin, J., y Quigley, A. (2021). \textit{Teacher feedback to improve pupil learning}. EEF.
    \item Fullan, M., y Langworthy, M. (2014). \textit{A rich seam: How new pedagogies find deep learning}. Pearson.
    \item Quiroga Aguilar, L., y Lara Órdenes, E. (2022). El aprendizaje profundo como herramienta para cambio en la visión de aprendizaje de una cultura escolar. \textit{REA}, 12(1).
    \item Restrepo Gómez, B. (2004). La investigación-acción educativa y la construcción de saber pedagógico. \textit{Educación y Educadores}, (7), 45--55.
    \item Rivas, A., André, F., y Delgado, L. E. (2017). \textit{50 innovaciones educativas para escuelas}. CIPPEC / Santillana.
    \item Sawyer, R. K. (2008). Optimizing learning: Implications of learning sciences research. OECD.
  \end{itemize}
\end{frame}

\end{document}
